% Track e4 -- The b-field between the constituents of a rung, up the ladder.
% Verification: investigations/emergence/verify/e4_bfield_ladder.py (claims e4-1a ... e4-5e).
% New references: sections/bib_e4.tex.
\section{\texorpdfstring{The $b$-field up the ladder}{The b-field up the ladder}}
\label{e4:sec:main}

Section~9 of \cite{RE} studied the field that acts between the constituents of one rung, the copy field $b=B\otimes\id_6$, and settled what it
does to rates and lines on two and three constituents. This section follows that field up the ladder, $n=2,3,4,\dots$, and reads it in
the Fock picture of \cref{e1:sec:fock,e2:sec:quanta}. The field always acts between the constituents of one rung $\Cl(3n,3n)$; going up the
ladder means increasing $n$.

The answers, with their classes, are these. (1) Every compact binding-compatible field is the second quantisation of a one-quantum
Hamiltonian. It conserves the number $N$ of quanta, moves quanta between the constituents, and its many-quantum lines are sums of
one-quantum lines. The copy boosts create pairs, one quantum in each of the two constituents they join (fixed; \cref{e4:sec:fock}). (2) On $n$
constituents with distinct rates and all couplings nonzero, the invariants of the copy field are the $|B_{kl}|$ and the signs of the
products of $B$ around the cycles of the complete graph. Exactly $(n-1)(n-2)/2$ of these signs are independent, and all of them can be taken
time-odd. The Jarlskog-type determinant $\det[H,K]$ has time parity $(-1)^n$, so it detects time reversal only at odd $n$ (fixed;
\cref{e4:sec:cycles}). (3) At $n=4$ the copy group is $SO(4)$, and each of its two $\su2$ factors has half-integer isospin exactly on the states
with an odd number of quanta. The copy inversion acts as $(-1)^N$. Above $n=4$ the same holds for the $\so4$ of any four constituents, with the
quanta in those four. At no rung does a spinor of the copy algebra occur (fixed; \cref{e4:sec:isospin}). (4) For rates of multiplicities $n_1+\dots+n_r=n$ the removable part of the copy $\so n$ is $\bigoplus_i\so{n_i}$, and
$\sum_{i<j}n_in_j$ directions are physical (fixed; the Higgs-like reading is interpretive; \cref{e4:sec:rates}). (5) For a rotor with rates
$(\kappa_k)$, a frame bracket between copies $i,j$ and $k$ is admitted only if $\kappa_k=\kappa_i+\kappa_j$. A harmonic ladder of rates
$\kappa_k=k\kappa_1$ therefore admits non-abelian, nilpotent frames. From $n=4$ on the Jacobi identity among the frame brackets is a
quadratic condition (fixed; the reading is constructed; \cref{e4:sec:rotor}). \Cref{e4:sec:answer} says what the field can and cannot do for the emergence of particles.
Identifiers e4-$n$ refer to the \texttt{PASS} lines of \texttt{verify/e4\_bfield\_ladder.py} (\cref{e4:tab:claims}).

\paragraph{Notation.} As in Section~9 of \cite{RE}: $V_n=\R^n\otimes\R^6$, $\eta_n$, $G_n=\id_n\otimes G$, $z=(z_1,\dots,z_n)\in\C^{3n}$ with
$z_k=\bB_k+i\bE_k$, copy rotations $C^{kl}=(E_{lk}-E_{kl})\otimes\id_6$, copy boosts $(E_{lk}-E_{kl})\otimes G$, copy field $b=B\otimes\id_6$ with
$B\in\so n$, rates $K=\operatorname{diag}(\kappa_1,\dots,\kappa_n)$, $\tilde K=K\otimes\id_3$, and, for a compact binding-compatible $b$ acting on $z$ as the
real antisymmetric $R$, the frame Hamiltonian $H=\tilde K-iR$ of Theorem~9.3 of \cite{RE}. $\sigma$ is the spin lift to
$S_n=(\R^8)^{\otimes n}$, with $[\sigma(X),v]=Xv$ for $v\in V_n$, and $\hat C^{kl}=\sigma(C^{kl})$. The $3n$ modes are $p=(k,i)$, $k\le n$, $i\le3$, with
$a^\pm_{k,i}=\half(E_{(k,i)}\pm E_{(k,i+3)})$ (Proposition~7.2 of \cite{RE}). $N_{k,i}=a^+_{k,i}a^-_{k,i}$, $N^{(k)}=\sum_iN_{k,i}$ is the number in
constituent $k$, $N_{[i]}=\sum_kN_{k,i}$ the number in direction $i$, and $N=\sum_kN^{(k)}$. For a complex $3n\times3n$ matrix $M$ we write
$\mathcal N(M)=\sum_{p,q}M_{pq}\,a^+_pa^-_q$. The first basis vector of the realisation is the Fock vacuum $|0\rangle$ (e4-1a).

\subsection{Quanta between constituents: transfer, not creation}
\label{e4:sec:fock}

\Cref{e1:thm:Nrule} shows, at every rung, that the compact binding-compatible fields are exactly the binding-compatible fields that
conserve $N$, and \cref{e2:prop:hop,e2:prop:copyboost} give the hop and the pair at $n=2$. The theorem below adds the second-quantised form,
which says how the quanta move, and the bookkeeping by constituent and by direction at $n=3$.

\begin{theorem}[The field between constituents, second quantised]
\label{e4:thm:fock}
\begin{enumerate}[label=(\roman*),leftmargin=2.2em]
\item Let $X$ be a paired rotation, acting on $z$ as the real antisymmetric $R\in\so{3n}$. Then $\sigma(X)=\mathcal N(R)$. Hence $\sigma(X)$ commutes
with $N$, and $\sigma(X)\,a^+_q|0\rangle=\sum_pR_{pq}\,a^+_p|0\rangle$.
\item For the copy field $b=B\otimes\id_6$ and rates $K$, the operator $\mathsf H=\sum_k\kappa_kN^{(k)}-i\sigma(b)$ on $S_n\otimes\C$ equals
$\mathcal N(H)$, $H=(K-iB)\otimes\id_3$. It is Hermitian and commutes with $N$. Its block on $\{N=1\}$, in the basis $a^+_p|0\rangle$, is $H$, and its
eigenvalues on $\{N=m\}$ are the sums of $m$ eigenvalues of $H$ with distinct indices.
\item At $n=3$, $\hat C^{kl}=\sum_i\bigl(a^+_{l,i}a^-_{k,i}-a^+_{k,i}a^-_{l,i}\bigr)$. It commutes with $N$, with $N^{(k)}+N^{(l)}$, with $N^{(m)}$ for
$m\notin\{k,l\}$ and with every $N_{[i]}$, and not with $N^{(k)}$. The three lifts close on $\so3$ as on vectors, $[\hat C^{12},\hat C^{13}]=\hat C^{23}$.
\item At $n=3$, $\sigma\bigl((E_{lk}-E_{kl})\otimes G\bigr)=-\sum_i\bigl(a^+_{k,i}a^+_{l,i}+a^-_{l,i}a^-_{k,i}\bigr)$. It joins only sectors whose $N$ differs
by $2$, in both directions, and it commutes with $N^{(k)}-N^{(l)}$ and with $N^{(m)}$, $m\notin\{k,l\}$: every pair it creates has one quantum in
constituent~$k$ and one in constituent~$l$.
\item At $n=3$ the quarter turn $U=e^{\pi\hat C^{12}/2}=\tfrac18\prod_a(\id+2t_a)$, with $t_a=-\half\eta_{aa}E_{(1,a)}E_{(2,a)}$ the six commuting terms of
$\hat C^{12}$, is orthogonal, commutes with $N$ and $N^{(3)}$, and satisfies $UN^{(1)}U^{-1}=N^{(2)}$. For $\kappa_1\neq\kappa_2$ it changes the rate
operator $\sum_k\kappa_kN^{(k)}$ and keeps its spectrum.
\end{enumerate}
\end{theorem}

\begin{proof}
(i) A paired rotation acts in the same way on the temporal directions and on their partners, so $X(2a^\pm_q)=\sum_pR_{pq}\,2a^\pm_p$, and
$[\sigma(X),a^\pm_q]=\sum_pR_{pq}a^\pm_p$. The CAR give $[\mathcal N(R),a^+_q]=\sum_pR_{pq}a^+_p$ and $[\mathcal N(R),a^-_q]=-\sum_pR_{qp}a^-_p$,
which is the same because $R^{\!\top}=-R$. So $\sigma(X)-\mathcal N(R)$ commutes with every generator, and $S_n$ is irreducible, so it is a multiple
of $\id$. Both operators are traceless: $\sigma(X)$ is a sum of products of two anticommuting generators, and $\Tr(a^+_pa^-_q)=0$ for $p\neq q$, while
$R_{pp}=0$. Hence they are equal. Each term $a^+_pa^-_q$ conserves $N$, and $a^-_q|0\rangle=0$ gives the action on one quantum. (ii)
$N^{(k)}=\mathcal N(E_{kk}\otimes\id_3)$, and by (i) $\sigma(b)=\mathcal N(B\otimes\id_3)$, so $\mathsf H=\mathcal N(H)$. $N^{(k)}$ is real diagonal and
$\sigma(b)$ real antisymmetric, so $\mathsf H$ is Hermitian. For $H=WDW^\dagger$ with $W$ unitary, $b^\pm_r=\sum_pW_{pr}^{(\pm)}a^\pm_p$ ($W^{(+)}=W$,
$W^{(-)}=\bar W$) satisfy the CAR and $\mathsf H=\sum_rd_r\,b^+_rb^-_r$; its eigenvectors are $b^+_{r_1}\cdots b^+_{r_m}|0\rangle$ with $r_1<\dots<r_m$. (iii) is
(i) for $R=(E_{lk}-E_{kl})\otimes\id_3$. Each term moves one quantum between $(k,i)$ and $(l,i)$. (iv) A copy boost anticommutes with $K_3$, so its
lift is a combination of $a^+a^+$ and $a^-a^-$ terms (\cref{e2:lem:N}(ii)). It maps the plane of mode $(k,i)$ to that of mode $(l,i)$, which fixes
the modes paired, and the coefficients are computed. (v) $t_a^2=-\tfrac14$, so
$e^{\theta\hat C^{12}}=\prod_a(\cos\tfrac\theta2+2\sin\tfrac\theta2\,t_a)$. At $\theta=\tfrac\pi2$, $\Ad U$ maps $E_{(1,a)}$ to $\pm E_{(2,a)}$ and $E_{(2,a)}$ to
$\mp E_{(1,a)}$, and fixes $E_{(3,a)}$. $N^{(k)}$ is a function of the products $E_{(k,i)}E_{(k,i+3)}$, so $UN^{(1)}U^{-1}=N^{(2)}$, and conjugation keeps
spectra.
\end{proof}

The theorem is fixed. Part (i) is the standard realisation of $\so{3n}$ on the exterior algebra by one-body operators
\cite{Chevalley54,Fock32}, inherited, and its realisation in the tower is checked exactly for all $15$ and $36$ basis elements at $n=2,3$
(e4-1b) and for the copy rotations at $n=4$ (e4-3a). Part (ii) is checked exactly on $\{N=1\}$ at $n=3$, rates $(1,2,4)$ and
$(B_{12},B_{23},B_{31})=(\tfrac12,\tfrac13,\tfrac15)$, and the additivity on $\{N=2\}$ and $\{N=3\}$ to $10^{-9}$ (e4-1e). Parts (iii) to (v) are exact
(e4-1c, e4-1d, e4-1f); in (iv) a random paired boost of $\so{9,\C}$ also changes $N$ by exactly $\pm2$. The external $i$ in $\mathsf H$ is constructed,
as in \cref{e2:prop:hop}(iii). Under $z\mapsto\sum_pz_pa^+_p|0\rangle$ the one-quantum sector of $e^{-i\mathsf Ht}$ is the frame flow $z(t)=e^{-iHt}z(0)$ of
Theorem~9.3(i) of \cite{RE}.

\begin{corollary}[Transfer, not creation]
\label{e4:cor:transfer}
Let $b$ be a compact element of $\Cent(G_n)\cap\so{3n,3n}$, the class of gauge-admissible fields between the constituents (Theorem~9.3(ii) of
\cite{RE}; Corollary~2.4 of \cite{ES17}).
\begin{enumerate}[label=(\roman*),leftmargin=2.2em]
\item $e^{t\sigma(b)}$ conserves $N$ and acts on $\{N=m\}$ as the $m$-th exterior power of $e^{tR}$. It moves quanta between modes and
constituents, and never creates or annihilates one.
\item The binding-compatible elements that change $N$ are the paired boosts, among them the copy boosts. Each of their terms creates or
annihilates one pair, and a copy boost puts one quantum in each of the two constituents it joins.
\item Let $\hat U$ lift a compact $\eta_n$-isometry commuting with $G_n$, which acts on $z$ as a unitary $W$. A change of frame by $\hat U$ takes the
rate operator $\mathcal N(\tilde K)$ to $\hat U\mathcal N(\tilde K)\hat U^{-1}=\mathcal N(W\tilde KW^{-1})$. This operator commutes with $N$ and has the same
spectrum on every sector $\{N=m\}$, the sums of $m$ rates. The numbers $N^{(k)}$ are not kept.
\end{enumerate}
\end{corollary}

\begin{proof}
(i) The proof of \cref{e4:thm:fock}(i) holds at every $n$, and $\mathcal N(R)$ is the derivation of the exterior algebra induced by $R$, so its
exponential is $\Lambda(e^{tR})$. (ii) is \cref{e1:thm:Nrule}(ii), with the pairing of \cref{e4:thm:fock}(iv) at $n=3$ and of
\cref{e2:prop:copyboost}(iii) at $n=2$. (iii) $\Ad\hat U$ acts on vectors as the isometry, hence on the $a^+_p$ as $W$, which gives the formula. A compact isometry commuting with $G_n$
commutes with $K_n=-\eta_nG_n$, so $\hat U$ commutes with $N$, and conjugation keeps spectra. This is the Fock form of Theorem~9.3(iii) of
\cite{RE}, and \cref{e4:thm:fock}(v) exhibits a frame change that moves $N^{(1)}$ to $N^{(2)}$.
\end{proof}

The corollary is fixed. It answers the first question precisely. The gauge-admissible field between constituents transfers quanta and never
creates them. It does not change a rate either (Theorem~9.3(iii) of \cite{RE}). Which constituent carries a quantum depends on the frame
when the copy rotation is gauged, as \cref{e2:prop:hop} found at $n=2$. How many quanta there are, and the multiset of their rates, does not.
Pair creation across constituents needs a copy boost, which is a metric connection and not a gauge field (Proposition~9.4 of \cite{RE}).
Because $\mathsf H$ is one-body, its many-quantum lines carry no interaction energy. The field mixes quanta; it does not bind them.

\subsection{\texorpdfstring{Time-odd invariants on $n$ constituents}{Time-odd invariants on n constituents}}
\label{e4:sec:cycles}

Theorem~9.5 of \cite{RE} found on three constituents one invariant that two do not carry, the sign of $J=B_{12}B_{23}B_{31}$, odd under interior
time reversal. For a cycle $\gamma=(k_1,\dots,k_L)$ of distinct copy indices write $\Pi_\gamma B=B_{k_1k_2}B_{k_2k_3}\cdots B_{k_Lk_1}$; reversing the
orientation multiplies it by $(-1)^L$.

\begin{theorem}[The copy field on $n$ constituents]
\label{e4:thm:cycles}
Let $n\ge3$, let the rates be distinct, and let $b=B\otimes\id_6$ with every $B_{kl}\neq0$.
\begin{enumerate}[label=(\roman*),leftmargin=2.2em]
\item The endomorphisms of $V_n$ commuting with $G_n$, the diagonal $\lG$ and $K\otimes G$ are $\sum_k(a_k+b_kG)$ on copy $k$, and the $\eta_n$-isometries
among them are the $2^n$ copy reflections $\operatorname{diag}(s_1,\dots,s_n)\otimes\id_6$, $s_k=\pm1$, under which $B_{kl}\mapsto s_ks_lB_{kl}$.
\item Two fields with the same $|B_{kl}|$ are related by a copy reflection iff $\Pi_\gamma B$ has the same sign for both on every cycle $\gamma$. The
signs of the $c_n=\tfrac12(n-1)(n-2)$ triangles $B_{1k}B_{kl}B_{l1}$, $2\le k<l\le n$, are independent and determine all the others. The
invariants of the field are the $\tfrac12n(n-1)$ magnitudes and these $c_n$ signs.
\item Interior time reversal, $(K,B)\mapsto(K,-B)$ (Theorem~9.5(iv) of \cite{RE}), multiplies $\Pi_\gamma B$ by $(-1)^L$. It flips all $c_n$ triangle signs,
and no copy reflection undoes it.
\item The triangle invariants are spectral. With $P_k$ the spectral projectors of $K$, polynomials in $K$, and $H=K-iB$ on the copy index,
$\tr(P_kHP_lHP_mH)=iB_{kl}B_{lm}B_{mk}$ for every triangle $(k,l,m)$.
\end{enumerate}
\end{theorem}

\begin{proof}
(i) The commutant of $G_n$ and the diagonal $\lG$ is $\{A\otimes\id+A'\otimes G\}$ with $A,A'\in\mathfrak{gl}(n)$ (Proposition~9.2 of \cite{RE}).
$[A\otimes\id+A'\otimes G,K\otimes G]=[A,K]\otimes G-[A',K]\otimes\id$, so $A$ and $A'$ commute with $K$ and are diagonal. The isometry condition per copy is
$a_kb_k=0$, $a_k^2-b_k^2=1$, as in the proof of Theorem~9.5(iii) of \cite{RE}. (ii) Encode the signs of $B_{kl}$, $k<l$, as $\sigma\in\Z_2^E$, $E$ the
edges of the complete graph $K_n$. A reflection adds $\delta s$, $(\delta s)_{kl}=s_k+s_l$, so the orbits are the cosets of $\operatorname{im}\delta$.
$\ker\delta$ is the constants because $K_n$ is connected, so $\rank\delta=n-1$ and there are $2^{|E|-n+1}=2^{c_n}$ orbits. The functions of $\sigma$ that are
constant on cosets are the characters $\sigma\mapsto\sum_{e\in\gamma}\sigma_e$ with $\gamma$ in the annihilator of $\operatorname{im}\delta$, the edge sets of even
degree at every vertex, which form the cycle space \cite{Diestel17}. Its dimension is $|E|-n+1=c_n$, and the fundamental cycles of the star at
vertex $1$, the triangles $(1,k,l)$, form a basis. A character pairs perfectly with the cosets, so the triangle signs separate the orbits. The
oriented product $\Pi_\gamma B$ differs from the unoriented one by a sign fixed by $\gamma$. (iii) $-B$ adds the all-ones vector to $\sigma$, and its
character on $\gamma$ is $L\bmod2$. A triangle is odd, so the all-ones vector is not in $\operatorname{im}\delta$. (iv) $P_kHP_l=H_{kl}E_{kl}$ and
$H_{kl}=-iB_{kl}$ for $k\neq l$, and $P_k=\prod_{j\neq k}(K-\kappa_j)/(\kappa_k-\kappa_j)$ for distinct rates.
\end{proof}

The theorem is fixed. Part (i) is exact at $n=4$, rates $(1,2,3,5)$ (e4-2a), and at $n=3$ it is Theorem~9.5(iii) of \cite{RE}. Parts (ii) and (iii)
are checked exhaustively on all sign patterns at $n=3,4,5,6$, with $c_n=1,3,6,10$ and $\rank\delta=n-1$ over $\Z_2$ (e4-2b, e4-2c). Part (iv) is
symbolic at $n=4$, where the product of the four triangle invariants is a product of squares, because every edge of $K_4$ lies on two
triangles (e4-2d). The count $c_n$ is the cycle rank of $K_n$, inherited from graph theory. It is also the number of rephasing-invariant phases of $n$
families of quarks \cite{KobayashiMaskawa73,BrancoLavouraSilva99} (inherited). Here the phases are fixed at $\pm\pi/2$ because the field is real
(Section~9.5 of \cite{RE}), so they reduce to signs. On $n$ constituents the copy field carries $c_n$ time-odd signs, its Wilson loops around the
triangles.

\begin{remark}[Couplings that vanish]
\label{e4:rem:graphs}
For a general coupling graph $\Gamma_B$, with $|E'|$ nonzero couplings and $c(\Gamma_B)$ connected components, the same argument gives
$2^{|E'|-n+c(\Gamma_B)}$ classes of sign patterns, so the number of independent cycle signs is the cycle rank $|E'|-n+c(\Gamma_B)$ (exhaustive over all
$64$ graphs at $n=4$, rank over $\Z_2$ for all graphs at $n=5,6$; e4-2f). A chain of constituents carries no sign. A ring of $n$ constituents
carries one sign, odd under time reversal iff $n$ is odd, in agreement with the bipartite criterion of Theorem~9.5(iv) of \cite{RE}. This is fixed.
\end{remark}

\begin{proposition}[The Jarlskog-type determinant up the ladder]
\label{e4:prop:jarlskog}
On the copy index $[H,K]=iD$ with $D_{kl}=(\kappa_k-\kappa_l)B_{kl}$, a real symmetric matrix with zero diagonal. Hence $\det[H,K]=i^n\det D$, and
$\det D$ is the sum over derangements $\pi$ of $(-1)^{n-\#\text{cycles}(\pi)}\prod_kD_{k\pi(k)}$. Under $B\mapsto-B$ it changes by $(-1)^n$.
\begin{enumerate}[label=(\roman*),leftmargin=2.2em]
\item $n=3$: $\det[H,K]=2iJ(\kappa_1-\kappa_2)(\kappa_1-\kappa_3)(\kappa_2-\kappa_3)$.
\item $n=4$: $\det[H,K]=D_{12}^2D_{34}^2+D_{13}^2D_{24}^2+D_{14}^2D_{23}^2-2(D_{12}D_{23}D_{34}D_{41}+D_{12}D_{24}D_{43}D_{31}+D_{13}D_{32}D_{24}D_{41})$. It is even
under time reversal and contains no triangle.
\item $n=5$: $\det[H,K]$ is odd under time reversal; every cycle cover contains a $5$-cycle or a triangle.
\end{enumerate}
\end{proposition}

\begin{proof}
$[K-iB,K]=-i[B,K]$ and $[B,K]_{kl}=-(\kappa_k-\kappa_l)B_{kl}$. $D$ is a product of two antisymmetric factors entrywise, hence symmetric. The Leibniz
expansion of a matrix with zero diagonal runs over derangements, and $-D$ multiplies each term by $(-1)^n$. At $n=4$ the derangements are the three
double transpositions and the six $4$-cycles. At $n=5$ they are the $24$ five-cycles and the $20$ products of a transposition and a $3$-cycle.
\end{proof}

The proposition is fixed. Part (i) restates Theorem~9.5(ii) of \cite{RE}, and the identities are symbolic at $n=3,4,5$, with the cycle-cover
sum checked at $n=4,5$ (e4-2e). The determinant measures cycle covers. A $2$-cycle contributes the magnitude $(\kappa_k-\kappa_l)^2B_{kl}^2$, and a
cycle of length $L\ge3$ contributes twice the product of the rate differences around it times $\Pi_\gamma B$. At $n=4$ the $4$-cycle products are
time-even: $B_{12}B_{23}B_{34}B_{41}=-B_{12}B_{23}B_{31}\cdot B_{13}B_{34}B_{41}/B_{13}^2$ is fixed by the product of two triangle signs. So
$\det[H,K]$ cannot tell $(K,B)$ from $(K,-B)$ at $n=4$, although the field is time-odd. The time-odd invariants there are the triangle traces of
\cref{e4:thm:cycles}(iv). Jarlskog's determinant \cite{Jarlskog85} measures time reversal on three constituents and, among the rungs, only at odd
$n$; the triangle traces measure it at every $n\ge3$.

\subsection{Copy isospin at four constituents}
\label{e4:sec:isospin}

On two and three constituents the copy isospin is integral (Proposition~9.6 of \cite{RE}), and Remark~9.8(d) of \cite{RE} notes that this does
not extend to four (d7-6d). We determine the copy isospin content of $S_4$, of dimension $4096$. Write
\[
I^\pm_1=\half(\hat C^{23}\pm\hat C^{14}),\qquad I^\pm_2=\half(\hat C^{31}\pm\hat C^{24}),\qquad I^\pm_3=\half(\hat C^{12}\pm\hat C^{34})
\]
(the self-dual and anti-self-dual combinations; $J^\pm_a$ in the script), and let $j_\pm(j_\pm+1)$ be the eigenvalues of the Casimirs
$-\sum_a(I^\pm_a)^2$.

\begin{theorem}[Copy isospin on four constituents]
\label{e4:thm:so4}
Let $n=4$.
\begin{enumerate}[label=(\roman*),leftmargin=2.2em]
\item Each $\hat C^{kl}$ equals $\sum_i(a^+_{l,i}a^-_{k,i}-a^+_{k,i}a^-_{l,i})$, commutes with every $N_{[i]}$, and is a sum of six commuting terms of square
$-\tfrac14$, so $e^{2\pi\hat C^{kl}}=+\id$. The six close on $\so4$, and the copy group acting on $S_4$ is $SO(4)$.
\item $[I^+_a,I^-_b]=0$ and $[I^\pm_1,I^\pm_2]=I^\pm_3$ cyclically, on vectors and on $S_4$. On vectors $4(I^\pm_a)^2=-\id$, so the copy index is the module
$(\half,\half)$. On $S_4$, $e^{2\pi I^+_3}=e^{2\pi I^-_3}=(-1)^N$. Hence $I^+_3=\half(\hat C^{12}+\hat C^{34})$ has eigenvalues $i\mu$ with $\mu$ half-integral
exactly on the sectors of odd $N$. In particular $e^{2\pi I^+_3}a^+_{1,1}|0\rangle=-a^+_{1,1}|0\rangle$.
\item As an $\so4$-module, $S_4\otimes\C\cong(\Lambda\C^4)^{\otimes3}$ with $\Lambda\C^4=2(0,0)\oplus2(\half,\half)\oplus(1,0)\oplus(0,1)$. The modules $(j_+,j_-)$ and their
multiplicities are those of \cref{e4:tab:so4}. Always $j_+-j_-\in\Z$. The modules with half-integral $j_+$ fill exactly the odd-$N$ sectors, $2048$
dimensions. The one-quantum sector is $(\half,\half)\otimes\R^3$.
\end{enumerate}
\end{theorem}

\begin{table}[htbp]
\centering
\footnotesize
\begin{tabular}{@{}lcccccccccc@{}}
\toprule
even $N$: $(j_+,j_-)$ & $(0,0)$ & $(1,0)$ & $(0,1)$ & $(1,1)$ & $(2,0)$ & $(0,2)$ & $(2,1)$ & $(1,2)$ & $(3,0)$ & $(0,3)$\\
multiplicity & $70$ & $84$ & $84$ & $90$ & $20$ & $20$ & $15$ & $15$ & $1$ & $1$\\
\midrule
odd $N$: $(j_+,j_-)$ & $(\half,\half)$ & $(\tfrac32,\half)$ & $(\half,\tfrac32)$ & $(\tfrac32,\tfrac32)$ & $(\tfrac52,\half)$ & $(\half,\tfrac52)$ & & & &\\
multiplicity & $140$ & $64$ & $64$ & $20$ & $6$ & $6$ & & & &\\
\bottomrule
\end{tabular}
\caption{The copy isospin content of $S_4\otimes\C$ under $\so4=\su2_+\oplus\su2_-$. Each half has dimension $2048$ (e4-3c).}
\label{e4:tab:so4}
\end{table}

\begin{proof}
(i) The one-body form is \cref{e4:thm:fock}(i), and the six terms are those of Proposition~9.6(i) of \cite{RE} with $|D|=6$. A module of $\so4$ on which
the $2\pi$ rotation in one plane acts as $\id$ integrates to $SO(4)$, not to its double cover. (ii) On vectors $2I^\pm_3=(C^{12}\pm C^{34})$ squares to
$-\id$, so $e^{2\pi I^\pm_3}=-\id$ on $V_4$. Its lift is the product of the twelve commuting factors $e^{2\pi u}$, where $u$ runs over the terms of
$I^\pm_3$, $u^2=-\tfrac1{16}$ and $e^{2\pi u}=4u$. That product is a multiple of the volume element of $\Cl(12,12)$, and the computation gives $(-1)^N$. On a
sector where $e^{2\pi I^+_3}=-\id$ the eigenvalues $i\mu$ of $I^+_3$ satisfy $e^{2\pi i\mu}=-1$. (iii) $S_4$ is the exterior algebra on the modes, and the
modes of direction $i$ span a copy $\C^4_i$ of the copy index, so $S_4\otimes\C=\Lambda(\C^4_1\oplus\C^4_2\oplus\C^4_3)=\bigotimes_i\Lambda\C^4_i$. By (i), $\so4$ acts
on each factor as the derivation induced by the vector module. $\Lambda^0=\Lambda^4=(0,0)$, $\Lambda^1\cong\Lambda^3=(\half,\half)$ and
$\Lambda^2=\Lambda^2_+\oplus\Lambda^2_-=(1,0)\oplus(0,1)$ \cite{Humphreys72}. The Clebsch--Gordan rule gives the table. Odd degree in $\Lambda\C^4$ means half-integral
$(j_+,j_-)$, and the parity of a product is the parity of the total degree, $(-1)^N$.
\end{proof}

The theorem is fixed. (i) and (ii) are exact on $S_4$, with sparse integer matrices (e4-3a, e4-3b), and $e^{2\pi\hat C^{12}}=+\id$ also on $S_5$, of dimension
$32768$ (e4-3a). (iii) is checked on the joint spectrum of the two Casimirs, block-diagonalised by the direction numbers $N_{[i]}$, which commute with
every copy rotation; the eigenvalues are integral to $10^{-8}$ and the multiplicities agree with the character (e4-3c). The decomposition rule is
inherited. The same rule at $n=3$, $\Lambda\C^3=2(j{=}0)\oplus2(j{=}1)$, gives the Casimir spectrum $j=0,1,2,3$ on $40,216,200,56$ dimensions of
Proposition~9.6(ii) of \cite{RE}, recomputed here (e4-3d, restated).

\begin{corollary}[No spinor of the copy algebra at any rung]
\label{e4:cor:nospinor}
For every $n$, $e^{2\pi\hat C^{kl}}=+\id$ on $S_n$, and $S_n\otimes\C\cong(\Lambda\C^n)^{\otimes3}$ as a module of the copy $\so n$. The copy algebra acts through
$SO(n)$, and for $n\ge3$ no spinor module of $\so n$ occurs. At $n=4$ the half-integral isospin is that of the vector $(\half,\half)$ and of its odd
exterior powers. At $n=6$, where $\so6\cong\su4$, the fundamental $\mathbf 4$ does not occur.
\end{corollary}

\begin{proof}
The six terms of $\hat C^{kl}$ commute and square to $-\tfrac14$ at every $n$, and the proof of \cref{e4:thm:so4}(iii) holds with $\C^n$ in place of $\C^4$.
Exterior powers of the vector module have integral weights in the orthonormal basis, and the spinor modules have half-integral ones
\cite{LawsonMichelsohn89}.
\end{proof}

The corollary is fixed, checked on $S_4$ and $S_5$ (e4-3a). It makes the doublets of Remark~9.8(d) of \cite{RE} precise. At $n=4$ a state carries
half-integral copy isospin under either $\su2$ factor iff it has an odd number of quanta; the copy inversion $-\id\in SO(4)$, the rotation through $\pi$
in the planes $(1,2)$ and $(3,4)$, acts as $(-1)^N$. This is possible without any double cover because $\so4$ is not simple and its vector is a
doublet of each factor. Up the ladder the copy algebra $\so n$ is simple again ($n\ge5$), so it has no $\su2$ ideal. Its subalgebra $\so4$ on any four
constituents still has doublets: on $S_5$, $e^{2\pi I^\pm_3}$ built on the constituents $(1,2,3,4)$ or $(2,3,4,5)$ is $(-1)^{N'}$, with $N'$ the number of
quanta in those four constituents (e4-3e). Half-integral copy isospin therefore first appears at the fourth rung and persists above it, under
the $\so4$ of four constituents. The doublets are frames and their odd products, not spinors of a spin cover.

\subsection{Rate patterns and the breaking of the copy symmetry}
\label{e4:sec:rates}

\begin{proposition}[Rate patterns]
\label{e4:prop:rates}
Let the rates take $r$ distinct values with multiplicities $n_1+\dots+n_r=n$.
\begin{enumerate}[label=(\roman*),leftmargin=2.2em]
\item The removable part of the copy $\so n$, the part commuting with $K\otimes G$, is $\bigoplus_i\so{n_i}$, block diagonal on the equal-rate
constituents, of dimension $\sum_i\half n_i(n_i-1)$. The non-removable directions number $\half n(n-1)-\sum_i\half n_i(n_i-1)=\sum_{i<j}n_in_j$. Within the
compact binding-compatible $\so{3n}$ the removable part is $\bigoplus_i\so{3n_i}$, and the non-removable directions number $9\sum_{i<j}n_in_j$.
\item The quadratic form $B\mapsto\|[K,B]\|^2$ on $\so n$ is diagonal in the basis $C^{kl}$ (on the copy index), with eigenvalue $2(\kappa_k-\kappa_l)^2$.
\end{enumerate}
\end{proposition}

\begin{proof}
(i) $[K,E_{lk}-E_{kl}]=(\kappa_l-\kappa_k)(E_{lk}+E_{kl})$, so $X\in\so n$ commutes with $K$ iff $X_{kl}=0$ whenever $\kappa_k\neq\kappa_l$, and likewise for $R$ and
$\tilde K$ on $\C^{3n}$. Removability is commutation with $K\otimes G$ (Theorem~9.3(v) of \cite{RE}). The count uses $n^2=\sum n_i^2+2\sum_{i<j}n_in_j$. (ii) The
matrices $E_{lk}+E_{kl}$, $k<l$, are orthogonal, of norm squared $2$.
\end{proof}

The proposition is fixed, exact for every partition of $n=2,\dots,5$ in $\so n$ and of $n=2,3,4$ in $\so{3n}$ (e4-4a), and at rates $(1,1,2,5)$, where the
form is $\operatorname{diag}(0,2,32,2,32,18)$ on $C^{12},C^{13},C^{14},C^{23},C^{24},C^{34}$ (e4-4b). At $n=3$ it restates the dimensions $36,18,9$ and $3,1,0$ of
Theorem~9.3(v) of \cite{RE}. With the factor $6$ of $\id_6$, (ii) is the norm $6\sum_{k\neq l}(\kappa_k-\kappa_l)^2B_{kl}^2$ of Remark~9.8(e) of \cite{RE}.

\begin{remark}[The breaking pattern]
\label{e4:rem:higgs}
The rate matrix is a symmetric tensor of the copy $\so n$, and its stabiliser is $\bigoplus_i\so{n_i}$; its stabiliser in $O(n)$ is
$O(n_1)\times\dots\times O(n_r)$, which for distinct rates is the group of copy reflections of \cref{e4:thm:cycles}(i). Read as in Remark~9.8(e) of
\cite{RE}, this is the pattern of a gauge symmetry broken by a symmetric-tensor vacuum value \cite{LiLF74}, with $\sum_{i<j}n_in_j$ massive directions
of masses proportional to $|\kappa_k-\kappa_l|$ \cite{EnglertBrout64,Higgs64,Kibble67}. Up the ladder, constituents with equal rates form a multiplet
whose internal copy rotations are removable, so which of them carries a quantum is a gauge choice. Constituents with different rates are
distinguishable, and the field between them is physical. The reading is interpretive, with the caveat of \cite{RE}: $b$ is a constant connection
without kinetic term, and the rates are inputs.
\end{remark}

\subsection{A ladder of rates admits a non-abelian frame}
\label{e4:sec:rotor}

Section~9.2 of \cite{RE} shows that a common binding action lets the rotor act by a non-central $A\in\uu n$. It also shows that the frame then
admits equivariant brackets $[X,Y]_k=\sum_{i,j}\Gamma^k_{ij}\beta_0(X_i,Y_j)$, with $\beta_0(x,y)=z(x)\times z(y)$ the complex cross product and
$\Gamma^k_{ij}=\Gamma^k_{ji}\in\C$, subject to $\Gamma^k_{ij}(\omega_k-\omega_i-\omega_j)=0$, where $A$ acts on copy $k$ as $i\omega_k$. For $A=K\otimes G$ the
$\omega_k$ are the rates. With a copy field, $A$ acts on $z$ as $iH$, and the condition applies to the eigenvalues of $H$, the lines.

\begin{proposition}[Rotor-compatible frame brackets]
\label{e4:prop:rotor}
\begin{enumerate}[label=(\roman*),leftmargin=2.2em]
\item The real bilinear maps $V\times V\to V$ equivariant under $\lG$ form the two-dimensional space spanned by $\beta_0$ and $G\beta_0$, and all of them
are antisymmetric. The equivariant antisymmetric brackets on $V_n$ are the maps above, a space of real dimension $n^2(n+1)=12,36,80$ for $n=2,3,4$.
\item $A=K\otimes G$ is a derivation of the bracket iff $\Gamma^k_{ij}=0$ unless $\kappa_k=\kappa_i+\kappa_j$. The rotor-compatible brackets have real dimension
$0$, $0$, $2$ for rates $(1,1)$, $(1,3)$, $(1,2)$; $0$, $4$, $0$, $4$ for $(1,1,1)$, $(1,2,3)$, $(1,3,5)$, $(1,2,4)$; and $8$ for $(1,2,3,4)$. For the
harmonic ladder $\kappa_k=k$ the dimension is $2\sum_{k\le n}\lfloor k/2\rfloor=2,4,8,12,18$ for $n=2,\dots,6$.
\item Rates $(1,2,3)$. The rotor-compatible brackets are $[X,Y]_2=a\,\beta_0(X_1,Y_1)$, $[X,Y]_3=b\,(\beta_0(X_1,Y_2)+\beta_0(X_2,Y_1))$, $(a,b)\in\C^2$. Every one of them
satisfies the Jacobi identity. For $ab\neq0$ the frame is a nilpotent Lie algebra $V_1\oplus V_2\oplus V_3$, graded by the rate and generated by $V_1$, with
lower central series of real dimensions $18,12,6,0$.
\item Rates $(1,2,4)$. The rotor-compatible brackets are $[X,Y]_2=a\,\beta_0(X_1,Y_1)$, $[X,Y]_3=b\,\beta_0(X_2,Y_2)$. For $X,Y\in V_1$ and $Z\in V_2$ the
Jacobiator is $ab\,\beta_0(\beta_0(X,Y),Z)\in V_3$, so the Jacobi identity holds iff $ab=0$.
\item Rates $(1,2,3,4)$, with $a=\Gamma^2_{11}$, $b=\Gamma^3_{12}$, $c=\Gamma^4_{13}$, $d=\Gamma^4_{22}$. For $X,Y\in V_1$ and $Z\in V_2$ the Jacobiator is
$(ad-bc)\,\beta_0(\beta_0(X,Y),Z)\in V_4$, and the Jacobi identity holds on all of $V_4$ iff $ad=bc$.
\end{enumerate}
\end{proposition}

\begin{proof}
(i) is an exact null-space computation on the $216$ coefficients of a bilinear map (e4-5a); $G$ acts as $i$, so $G\beta_0=i\beta_0$, and the maps are complex
bilinear. A bracket is a choice of $\Gamma^k_{ij}$ for each $k$ and each unordered pair $\{i,j\}$, $n\cdot\half n(n+1)$ complex numbers. (ii) On $z$,
$A$ acts as $i\kappa_k$ on copy $k$, and the bracket is complex bilinear, so $A[X,Y]_k-[AX,Y]_k-[X,AY]_k=i(\kappa_k-\kappa_i-\kappa_j)\Gamma^k_{ij}\beta_0(X_i,Y_j)$. For
$\kappa_k=k$ the pairs $i\le j$ with $i+j=k$ number $\lfloor k/2\rfloor$. (iii) The only composite is $[[X,Y]_2,Z]_3$, so the Jacobiator is
$ab\sum_{\rm cyc}\beta_0(\beta_0(X_1,Y_1),Z_1)$, which vanishes by the Jacobi identity of the complex cross product, the bracket of $\so{3,\C}$.
$\beta_0(V_1,V_1)=V_2$ and $\beta_0(V_1,V_2)=V_3$ for $ab\neq0$. The Jacobiator is trilinear and alternating, and inputs of rates $r_1,r_2,r_3$ give an output of rate $r_1+r_2+r_3$, so each of its components is fixed
by inputs whose rates add to the target rate. (iv) The only composite is $[[X,Y]_2,Z]_3=ab\,\beta_0(\beta_0(X_1,Y_1),Z_2)$, with inputs of rates $1,1,2$. For $X,Y\in V_1$ and $Z\in V_2$
the other two cyclic terms vanish, because $V_1\times V_2$ brackets to zero. For $ab=0$ there is no composite. (v) The $V_3$-component is as in (iii)
and vanishes. For $X,Y\in V_1$ and $Z\in V_2$, $[[X,Y],Z]_4=ad\,\beta_0(\beta_0(X,Y),Z)$, while
$[[Y,Z],X]_4+[[Z,X],Y]_4=bc\,(\beta_0(\beta_0(Y,Z),X)+\beta_0(\beta_0(Z,X),Y))=-bc\,\beta_0(\beta_0(X,Y),Z)$ by the Jacobi identity of $\beta_0$.
\end{proof}

The proposition is fixed. (i) restates the space of Section~9.2 of \cite{RE} at $n=2$ (e4-5a). (ii) is exact, block by block, for the listed rates and
for the harmonic ladder up to $n=4$; for $n=5,6$ the count is the number of resonant blocks (e4-5b). At $n=2$ it restates the bracket of
$\operatorname{diag}(1,2)\otimes G$ of Section~9.2 of \cite{RE}. (iii) to (v) are checked on random samples to $10^{-10}$: antisymmetry,
$\lG$-equivariance, the rotor as a derivation and the Jacobi identity; at $(1,2,4)$ and $(1,2,3,4)$ the samples off the stated conditions fail (e4-5c,
e4-5d, e4-5e). That every rotor-compatible frame is nilpotent when all rates are positive is inherited: $A$ is then a nonsingular derivation, and a
Lie algebra with a nonsingular derivation is nilpotent \cite{Jacobson55}.

The reading is constructed. A harmonic ladder of rates admits a non-abelian frame, from $n=3$ on and already for $\operatorname{diag}(1,2)$ at $n=2$. Its
bracket joins two frames whose rates add into the frame of the summed rate, so the frame algebra is graded by rate. With three constituents
on the harmonic ladder the Lie frame brackets form the linear space $\C^2$. With four the Jacobi identity cuts out the cone $ad=bc$ in $\C^4$, and
on a longer ladder the same condition recurs, because the $V_4$-component of the Jacobiator involves only the blocks with target rate at most
$4$. The resonant pattern $(1,2,4)$ admits a non-abelian frame only with one of its two blocks, $ab=0$. Nothing selects such a rotor. The paper's closure takes $A=G_n$ and an abelian frame
(Theorem~17.7 and Proposition~17.18 of \cite{ES17}). The non-central rotor, its rate ladder and the bracket are admitted, not selected. The
resonance condition $\omega_k=\omega_i+\omega_j$ is that of three-wave mixing and sum-frequency generation \cite{Armstrong62}, a parallel only
(interpretive). Read with $mc^2=\hbar\kappa$, the bracket composes frames whose rest energies add (interpretive).

\subsection{Answer: what the field between constituents can and cannot do}
\label{e4:sec:answer}

\begin{remark}[What the field does up the ladder]
\label{e4:rem:answer}
\begin{enumerate}[label=(\alph*),leftmargin=2.2em]
\item \emph{It moves quanta and keeps their number} (fixed). A gauge-admissible field between the constituents acts as a one-body operator.
It hops quanta between constituents and mixes them, and its $N$-quantum lines are sums of one-quantum lines (\cref{e4:thm:fock} and
\cref{e4:cor:transfer}).
\item \emph{It does not take one to two} (fixed). It creates no quantum and changes no rate. Pair emergence across constituents needs a copy boost,
which is not a gauge field (\cref{e4:cor:transfer}(ii); Proposition~9.4 of \cite{RE}). Because it is one-body, it binds no quanta into composites.
If two particles are to emerge from one, the process lies outside the compact field between constituents.
\item \emph{It carries time-odd invariants that grow with the rung} (fixed). From $n=3$ on it carries $\tfrac12(n-1)(n-2)$ independent triangle signs, each odd
under interior time reversal. The triangle traces measure them at every $n$, Jarlskog's determinant only at odd $n$
(\cref{e4:thm:cycles,e4:prop:jarlskog}).
\item \emph{It acts on doublets from the fourth rung on} (fixed). At $n=4$ the copy isospin is half-integral exactly on odd $N$, and above it under the
$\so4$ of any four constituents, on an odd number of quanta in them. At every rung the copy algebra acts through $SO(n)$, so no spinor of it
occurs (\cref{e4:thm:so4,e4:cor:nospinor}).
\item \emph{It is physical exactly between unequal rates} (fixed). The physical directions number $\sum_{i<j}n_in_j$; equal-rate constituents stay
gauge-equivalent (\cref{e4:prop:rates}). The reading as symmetry breaking by the rate matrix is interpretive (\cref{e4:rem:higgs}).
\item \emph{It admits non-abelian frames on a harmonic ladder of rates} (constructed). The rotor must be non-central and its rates resonant, and
nothing in the closure selects either (\cref{e4:prop:rotor}).
\end{enumerate}
Up the ladder, then, the field between constituents distinguishes constituents by their rates, carries more time-odd structure, and from
$n=4$ on acts on doublets. It never changes the number of quanta that the rung carries. In the sense of \cref{e1:sec:fock}, a rung carries $N$
quanta, and the compact $b$-field keeps $N$ fixed. It decides where the quanta are, not how many there are. This reading is interpretive, drawn
from the fixed statements above.
\end{remark}

\begin{table}[htbp]
\centering
\footnotesize
\begin{tabular}{@{}>{\raggedright\arraybackslash}p{0.4cm}>{\raggedright\arraybackslash}p{2.6cm}>{\raggedright\arraybackslash}p{2.0cm}>{\raggedright\arraybackslash}p{1.8cm}>{\raggedright\arraybackslash}p{3.8cm}>{\raggedright\arraybackslash}p{2.3cm}@{}}
\toprule
$n$ & copy algebra, group on $S_n$ & physical directions, distinct rates & independent signs $c_n$ & half-integral copy isospin & rotor brackets, rates $1,\dots,n$\\
\midrule
$2$ & $\so2$, $SO(2)$ & $1$ & $0$ & no; charges $i\{-3,\dots,3\}$ & $2$\\
$3$ & $\so3$, $SO(3)$ & $3$ & $1$ & no; $j=0,\dots,3$ & $4$\\
$4$ & $\su2\oplus\su2$, $SO(4)$ & $6$ & $3$ & yes, exactly on odd $N$ & $8$, Jacobi: $ad=bc$\\
$5$ & $\so5$, $SO(5)$ & $10$ & $6$ & under the $\so4$ of four constituents; no spinor of $\so5$ & $12$\\
$6$ & $\so6\cong\su4$, $SO(6)$ & $15$ & $10$ & as for $n=5$; no $\mathbf 4$ of $\su4$ & $18$\\
\bottomrule
\end{tabular}
\caption{The field between constituents, rung by rung. The $n=2,3$ isospin entries restate Proposition~9.6 of \cite{RE}. The others follow from
\cref{e4:thm:cycles,e4:thm:so4,e4:cor:nospinor,e4:prop:rates,e4:prop:rotor}; the rotor column counts real dimensions.}
\label{e4:tab:ladder}
\end{table}

\begin{table}[htbp]
\centering
\footnotesize
\begin{tabular}{@{}>{\raggedright\arraybackslash}p{8.0cm}>{\raggedright\arraybackslash}p{4.8cm}>{\raggedright\arraybackslash}p{2.0cm}@{}}
\toprule
statement & class & checks\\
\midrule
\cref{e4:thm:fock}(i): paired rotations are one-body operators $\mathcal N(R)$, conserve $N$ & fixed (structure inherited, realisation fixed) & e4-1a, e4-1b\\
\cref{e4:thm:fock}(ii): $\mathsf H=\mathcal N(H)$; $N=1$ block is $H$; additive lines & fixed; external $i$ constructed & e4-1e\\
\cref{e4:thm:fock}(iii), (iv): copy rotations and copy boosts at $n=3$ by constituent and direction & fixed & e4-1c, e4-1d\\
\cref{e4:thm:fock}(v): quarter turn moves $N^{(1)}$ to $N^{(2)}$, keeps the rate spectrum & fixed & e4-1f\\
\cref{e4:cor:transfer}: transfer, not creation; rates gauge invariant & fixed (partly restated: \cref{e1:thm:Nrule}; RE Thm~9.3(iii)) & e4-1b, e4-1d, e4-3a\\
\cref{e4:thm:cycles}(i): residual group is the copy reflections & fixed (restated at $n=3$: RE Thm~9.5(iii)) & e4-2a\\
\cref{e4:thm:cycles}(ii), (iii): invariants are $|B_{kl}|$ and $c_n$ triangle signs; all time-odd & fixed (cycle space inherited) & e4-2b, e4-2c\\
\cref{e4:thm:cycles}(iv): triangle traces & fixed & e4-2d\\
\cref{e4:rem:graphs}: cycle rank of the coupling graph & fixed & e4-2f\\
\cref{e4:prop:jarlskog}: $\det[H,K]=i^n\det D$, parity $(-1)^n$, $n=3,4,5$ & fixed ($n=3$ restated: RE Thm~9.5(ii)) & e4-2e\\
parallel with the $(n-1)(n-2)/2$ phases of $n$ quark families & inherited & ---\\
\cref{e4:thm:so4}(i), (ii): $SO(4)$, $(\half,\half)$, $e^{2\pi I^\pm_3}=(-1)^N$ & fixed & e4-3a, e4-3b\\
\cref{e4:thm:so4}(iii), \cref{e4:tab:so4}: $(\Lambda\C^4)^{\otimes3}$, doublets on odd $N$ & fixed (decomposition rule inherited) & e4-3c\\
$n=3$ Casimir spectrum $40,216,200,56$ & restated (RE Prop.~9.6(ii)) & e4-3d\\
\cref{e4:cor:nospinor}: no spinor of the copy algebra at any rung; doublets of the $\so4$ of four constituents at $n=5$ & fixed & e4-3a, e4-3e\\
\cref{e4:prop:rates}: removable $\bigoplus\so{n_i}$, $\sum_{i<j}n_in_j$ physical; mass form & fixed ($n=3$ restated: RE Thm~9.3(v)) & e4-4a, e4-4b\\
\cref{e4:rem:higgs}: breaking by the rate matrix & interpretive & ---\\
\cref{e4:prop:rotor}(i), (ii): equivariant brackets, resonance condition, dimensions & fixed ($n=2$ restated: RE Sec.~9.2) & e4-5a, e4-5b\\
\cref{e4:prop:rotor}(iii)--(v): Jacobi at $(1,2,3)$, $(1,2,4)$, $(1,2,3,4)$; nilpotency & fixed; nilpotency inherited & e4-5c, e4-5d, e4-5e\\
\cref{e4:sec:rotor}: a harmonic ladder admits a non-abelian frame & constructed; three-wave parallel interpretive & e4-5b, e4-5c\\
\cref{e4:rem:answer}, \cref{e4:tab:ladder}: what the field can and cannot do & (a)--(e) fixed, (f) constructed; summary interpretive & ---\\
\bottomrule
\end{tabular}
\caption{Claims of \cref{e4:sec:main}, their classes and their checks, which are \texttt{PASS} lines of \texttt{verify/e4\_bfield\_ladder.py}.}
\label{e4:tab:claims}
\end{table}
